\chapter{Background and State of the Art}
\label{ch:background}

\section{Evolution of Natural Language Processing}
Natural Language Processing (NLP) has undergone several revolutionary transformations over the past few decades. Early NLP systems relied on symbolic logic and heavily hand-crafted grammatical rules. These systems were brittle and struggled with the inherent ambiguity of human language.

The introduction of statistical machine learning, followed by the deep learning revolution, shifted the paradigm towards neural networks. Recurrent Neural Networks (RNNs) and Long Short-Term Memory (LSTM) networks became the de facto standard for sequence modeling. These architectures processed data sequentially, maintaining a hidden state that captured historical context. However, the sequential nature of RNNs presented two fundamental flaws: it precluded parallelization across the sequence length, making training on large datasets exceedingly slow, and it suffered from the vanishing gradient problem, wherein the network struggled to retain information over long contexts.

\section{The Transformer Architecture}
In 2017, Vaswani et al. \cite{vaswani2017attention} introduced the Transformer architecture, a purely attention-based model that completely discarded recurrence and convolutions. By processing all tokens in a sequence simultaneously, the Transformer unlocked massive parallelization capabilities, perfectly suited for modern GPUs.

\subsection{Self-Attention Mechanism}
At the heart of the Transformer is the Scaled Dot-Product Attention mechanism. For a given input sequence, the model generates three distinct representations for each token: a Query ($Q$), a Key ($K$), and a Value ($V$), typically via linear projections.

The attention score between any two tokens is computed as the dot product of their respective Query and Key vectors. To prevent the dot products from growing excessively large—which would push the subsequent softmax function into regions with extremely small gradients—the scores are scaled down by the square root of the hidden dimension size ($d_k$).

The mathematical formulation is given by:
\begin{equation}
\text{Attention}(Q, K, V) = \text{softmax}\left(\frac{QK^T}{\sqrt{d_k}}\right)V
\end{equation}

This operation effectively allows every token in the sequence to dynamically determine how much "attention" it should pay to every other token, capturing complex syntactic and semantic dependencies regardless of their distance in the text.

\subsection{Multi-Head Attention}
Instead of computing a single attention function, the Transformer employs Multi-Head Attention (MHA). The $Q$, $K$, and $V$ matrices are linearly projected $h$ times into lower-dimensional subspaces. The attention function is performed in parallel on each of these projected versions, and the resulting vectors are concatenated and linearly projected once more. This allows the model to jointly attend to information from different representation subspaces at different positions.

\subsection{Position-wise Feed-Forward Networks}
Following the attention sub-layer, each block in the Transformer applies a Position-wise Feed-Forward Network (FFN). This consists of two linear transformations with a non-linear activation function (such as ReLU or GeLU) in between:
\begin{equation}
\text{FFN}(x) = \max(0, xW_1 + b_1)W_2 + b_2
\end{equation}
While the attention mechanism aggregates information across different tokens, the FFN processes the aggregated information independently for each token position.

\section{Model Compression and DistilBERT}
While highly effective, Transformers suffer from quadratic scaling. The attention matrix $QK^T$ requires $O(N^2)$ computations and memory, where $N$ is the sequence length. Consequently, models like BERT (Bidirectional Encoder Representations from Transformers) are characterized by massive parameter counts—110 million for BERT-base and 340 million for BERT-large.

Deploying such monolithic models on resource-constrained edge devices is unfeasible. To address this, various model compression techniques have been developed, including pruning (removing redundant weights), quantization (reducing numerical precision), and knowledge distillation.

\subsection{Knowledge Distillation}
Knowledge Distillation is a compression technique in which a compact "student" model is trained to reproduce the behavior of a larger, pre-trained "teacher" model. Instead of training the student solely on the hard labels (e.g., classifying an image as a dog or a cat), the student is trained on the "soft probabilities" output by the teacher.

Sanh et al. \cite{sanh2019distilbert} applied this technique to BERT to create DistilBERT. By halving the number of layers from 12 to 6 and utilizing a specialized loss function combining language modeling, distillation, and cosine-distance losses, DistilBERT achieves a 40\% reduction in parameters (down to 66 million) and is 60\% faster during inference, while retaining 97\% of the original model's natural language understanding capabilities.

Despite this architectural compression, executing DistilBERT on edge devices via standard Python-centric deep learning frameworks like PyTorch or TensorFlow still incurs massive runtime overhead, necessitating compiler-level interventions.

\section{The Evolution of Compilers}
To understand the necessity of MLIR, one must first understand the architecture of traditional compilers.

\subsection{The Three-Phase Compiler Design}
Modern compilers, such as GCC and LLVM, are structured into three distinct phases:
\begin{enumerate}
    \item \textbf{Frontend:} Parses the source code (e.g., C, C++, Rust), performs lexical and semantic analysis, and translates the code into a generic Intermediate Representation (IR).
    \item \textbf{Optimizer (Middle-end):} Analyzes the IR and applies target-independent transformations to improve performance, such as dead code elimination, loop unrolling, and constant propagation.
    \item \textbf{Backend:} Translates the optimized IR into target-specific machine code (e.g., x86, ARM), performing instruction selection and register allocation.
\end{enumerate}

\subsection{LLVM IR and its Limitations for Machine Learning}
The LLVM compiler infrastructure revolutionized the industry by defining a strict, strongly-typed, SSA-based (Static Single Assignment) Intermediate Representation. LLVM IR is an excellent representation for scalar operations, memory addressing, and low-level control flow.

However, LLVM IR is fundamentally scalar. When a high-level Machine Learning framework attempts to lower a concept like a "Matrix Multiplication" to LLVM IR, the concept is destroyed, flattened into thousands of primitive \texttt{fadd} and \texttt{fmul} instructions hidden within complex loop nests. Once flattened, it becomes mathematically impossible for the compiler to look at the IR and deduce, "This is an attention mechanism; I should fuse it."

\section{The MLIR Framework}
The Multi-Level Intermediate Representation (MLIR) \cite{lattner2020mlir} was engineered to bridge the massive semantic gap between high-level Python ML frameworks and low-level machine code compilers. 

\subsection{The Concept of Dialects}
Instead of forcing all programs into a single IR, MLIR acts as a meta-compiler framework. It allows developers to define custom "dialects"—modular namespaces containing specific operations, types, and attributes tailored to a particular domain.

A compilation pipeline in MLIR does not lower code directly to machine code. Instead, it progressively "lowers" code through multiple dialects, applying the most appropriate optimizations at the exact level of abstraction where they make the most sense. 

For example, our pipeline utilizes:
\begin{itemize}
    \item \textbf{The \texttt{dbert} Dialect:} A custom, high-level dialect where a neural network layer is represented as a single instruction. At this level, it is trivial to perform graph topology optimizations like Attention Fusion.
    \item \textbf{The \texttt{linalg} (Linear Algebra) Dialect:} A mid-level dialect that abstracts loop iterations over structured data. At this level, loop tiling and memory layout transformations are performed.
    \item \textbf{The \texttt{vector} Dialect:} A lower-level dialect that models SIMD hardware concepts. At this level, operations are mapped to 128-bit vector chunks.
    \item \textbf{The \texttt{llvm} Dialect:} The lowest MLIR dialect, mapping 1:1 with LLVM IR, allowing the compiler to hand the IR over to the LLVM backend for final register allocation and machine code emission.
\end{itemize}

\section{The Mathematics of Quantization}
Quantization is critical for edge inference. Neural networks are traditionally trained using 32-bit floating-point arithmetic (FP32), which offers immense dynamic range. However, research has shown that the deep layers of neural networks are highly resilient to noise and loss of precision \cite{jacob2018quantization}.

Quantization maps continuous real numbers into a discrete set of integer values. 

\subsection{Affine Quantization Mapping}
The general formula for Affine (Asymmetric) Quantization is:
\begin{equation}
q = \text{round}\left(\frac{r}{S}\right) + Z
\end{equation}
where:
\begin{itemize}
    \item $r$ is the original real-valued floating-point number.
    \item $q$ is the resulting quantized integer (e.g., an 8-bit integer in the range [-128, 127] or [0, 255]).
    \item $S$ is the Scale factor, a positive 32-bit floating-point number representing the step size.
    \item $Z$ is the Zero-point, an integer representing the value that the real number $0.0$ maps to in the quantized space.
\end{itemize}

\subsection{Symmetric Quantization}
In this report, we utilize Symmetric Quantization for the parameter weights and activations. In symmetric quantization, the zero-point $Z$ is strictly constrained to 0. This constraint drastically simplifies the dot-product computations at runtime, as the CPU does not need to compute cross-multiplication offset terms.

For symmetric 8-bit signed quantization (INT8), the range is $[-127, 127]$. Given a calibrated minimum value ($r_{min}$) and maximum value ($r_{max}$) from an FP32 tensor, the maximum absolute range is defined as:
\begin{equation}
r_{max\_abs} = \max(|r_{min}|, |r_{max}|)
\end{equation}
The scale factor $S$ is then derived as:
\begin{equation}
S = \frac{r_{max\_abs}}{127.0}
\end{equation}

By transforming $O(N^3)$ matrix multiplications from FP32 to INT8, we achieve a $4\times$ reduction in memory bandwidth and a $4\times$ increase in L1 cache capacity utilization.

\section{Hardware Architecture: ARM Cortex-A72}
To effectively optimize software, the compiler must intimately understand the hardware. The target environment is the Broadcom BCM2711 System-on-Chip (SoC) housed within the Raspberry Pi 4, featuring a quad-core ARM Cortex-A72 processor.

\subsection{The Memory Hierarchy and Cache Thrashing}
Modern processors execute instructions exponentially faster than main memory (RAM) can supply data. To bridge this gap, CPUs employ a hierarchy of SRAM caches.
The Cortex-A72 features:
\begin{itemize}
    \item A \textbf{32KB L1 Data Cache} per core (extremely fast, 1-3 cycle latency).
    \item A \textbf{1MB L2 Shared Cache} (moderately fast, ~10 cycle latency).
    \item \textbf{LPDDR4 DRAM} (extremely slow, ~100+ cycle latency).
\end{itemize}

When calculating a matrix multiplication of size $M \times K$ by $K \times N$, the CPU must fetch row and column vectors. If the matrices are larger than 32KB, the inner loops will constantly evict data from the L1 cache, forcing the CPU to fetch from RAM. This is known as "cache thrashing." For DistilBERT, an attention probability matrix of size $128 \times 128$ in FP32 consumes exactly 64KB, severely thrashing the 32KB L1 cache.

\subsection{Advanced SIMD (NEON)}
The Cortex-A72 implements the ARMv8-A architecture, which includes the Advanced SIMD extension known as NEON. NEON features thirty-two 128-bit wide vector registers. 
Standard scalar execution processes one data element per instruction. NEON allows Single Instruction, Multiple Data (SIMD) execution. A single 128-bit register can be partitioned to hold four 32-bit floats, or sixteen 8-bit integers. 

By heavily relying on MLIR's Vector dialect, our compiler explicitly structures the computation loops to pack sixteen quantized values into these registers, allowing the processor to execute 16 multiplications in a single clock cycle.


